\label{sec:embedding-task}
Описанный в §1.2 пайплайн выдаёт для каждого изображения набор диаграмм персистентности --- по одной на каждое направление PHT-фильтрации, плюс, при желании, отдельную диаграмму подуровневой фильтрации. Чтобы использовать эти диаграммы в задачах классификации и регрессии, их нужно превратить в признаки, совместимые с алгоритмами машинного обучения. Эта и составляет \emph{задачу эмбеддинга} диаграмм персистентности.

\paragraph{Почему задача нетривиальна.}
Диаграмма персистентности $D$ --- это \emph{мультимножество} точек переменной мощности в $\mathbb{R}^2 \cup (\mathbb{R} \times \{\infty\})$. У такого объекта есть три свойства, затрудняющих прямое использование стандартных моделей:
\begin{enumerate}
\item \emph{Переменный размер}: разные изображения порождают разное число точек в диаграмме, тогда как большинство ML-методов требует входа фиксированной размерности.
\item \emph{Перестановочная инвариантность}: диаграмма как множество не имеет естественного порядка, и любой эмбеддинг $\phi(D)$ должен быть инвариантен к перестановке точек.
\item \emph{Невозможность изометрического вложения в гильбертово пространство}: пространство диаграмм с метрикой Вассерштейна, как показано в работах~\cite{mileykoProbabilityMeasuresSpace2011}, не вкладывается изометрически ни в какое гильбертово пространство. Это означает, что прямое использование диаграммы как точки в евклидовом пространстве не сохраняет геометрию пространства диаграмм; искажение неизбежно.
\end{enumerate}

\paragraph{Формальная постановка.}
Будем искать отображение $\phi: \mathcal{D} \to \mathbb{R}^n$ (или в гильбертово пространство $\mathcal{H}$), где $\mathcal{D}$ --- пространство диаграмм персистентности, удовлетворяющее набору желательных свойств:
\begin{itemize}
\item[(i)] \textbf{Фиксированная размерность.} Выход $\phi(D)$ --- вектор фиксированной длины, не зависящей от числа точек в $D$.
\item[(ii)] \textbf{Перестановочная инвариантность.} $\phi(D) = \phi(\pi \cdot D)$ для любой перестановки точек $\pi$ диаграммы $D$.
\item[(iii)] \textbf{Устойчивость.} Небольшое возмущение входа (в смысле метрики Вассерштейна на $\mathcal{D}$) приводит к небольшому изменению $\phi(D)$.
\item[(iv)] \textbf{Дифференцируемость} по обучаемым параметрам $\phi_\theta$, что позволяет интегрировать эмбеддинг в end-to-end обучаемые pipelines.
\item[(v)] \textbf{Достаточная выразительность} для конкретной задачи --- т.\,е. способность аппроксимировать произвольную непрерывную функцию на $\mathcal{D}$ (универсальная аппроксимация).
\end{itemize}

\paragraph{Два класса методов.}
В литературе сложилось два больших семейства подходов к эмбеддингу диаграмм:
\begin{itemize}
\item \emph{Предписанные (prescribed) представления} --- $\phi$ задаётся аналитически и не имеет обучаемых параметров. Такие методы, как правило, обладают доказанной устойчивостью и легко интерпретируемы, но не настраиваются под конкретную задачу и теряют информацию в шагах дискретизации и при выборе весовой функции. Типичные представители: persistence image~\cite{adamsPersistenceImages2017}, persistence landscape~\cite{bubenikStatisticalTopologicalData2015}, persistence silhouette~\cite{chazalStochasticConvergencePersistence2014}.
\item \emph{Обучаемые (learned) представления} --- $\phi_\theta$ параметризовано нейронной сетью, обученной совместно с downstream-задачей. Сюда относятся DeepSets~\cite{zaheerDeepSets2017}, PersLay и Persformer~\cite{reinauerPersformerTransformerArchitecture2021}. Они более выразительны и адаптивны, но требуют значительно больше данных и вычислений.
\end{itemize}

\paragraph{Особенности нашей постановки.}
В нашей работе вход усложняется тем, что для одного изображения формируется набор диаграмм по 16 направлениям PHT (а иногда плюс подуровневая фильтрация). С точки зрения эмбеддинга это означает, что:
\begin{itemize}
\item каждая точка приходит с метаинформацией о направлении $\alpha$ и о типе фильтрации, и эмбеддинг должен уметь использовать эту информацию;
\item общее число точек на одно изображение становится большим (сотни и тысячи), что вносит требования по масштабируемости в число точек;
\item инвариантность нужна именно по перестановкам внутри объединённого мультимножества, не по полностью независимым диаграммам.
\end{itemize}
Эти соображения и формируют требования к предлагаемой в §3 архитектуре. В §2 мы сначала разберём, как существующие методы эмбеддинга решают перечисленные критерии (i)--(v), а затем обсудим, где они оказываются недостаточны для нашего сценария.
